\documentclass[11pt]{article}

% --- Paquets de base ---
\usepackage[utf8]{inputenc}
\usepackage[T1]{fontenc}
\usepackage[french]{babel}
\usepackage{amsmath, amsfonts, amssymb}
\usepackage{geometry}
\geometry{a4paper, margin=1in}

% --- Mise en page et Style ---
\usepackage{titlesec}
\usepackage{enumitem}
\usepackage{booktabs}
\usepackage{longtable}
\usepackage{xcolor}
\usepackage{tcolorbox}
\usepackage{tikz}
\usepackage{hyperref}
\usepackage{caption}
\usepackage{graphicx}
\usetikzlibrary{shapes.geometric, arrows.meta, positioning, calc}

% --- Macros Mathématiques ---
\newcommand{\E}{\mathbb{E}}
\newcommand{\Prob}{\mathbb{P}}
\newcommand{\R}{\mathbb{R}}
\newcommand{\N}{\mathbb{N}}
\newcommand{\Var}{\mathrm{Var}}
\newcommand{\Cov}{\mathrm{Cov}}
\newcommand{\calN}{\mathcal{N}}
\newcommand{\calF}{\mathcal{F}}
\newcommand{\ind}{\mathbf{1}}

% --- Boîtes de définition ---
\tcbuselibrary{skins, breakable}
\newtcolorbox{definitionbox}[1][]{
  colback=gray!8, colframe=black!60,
  fonttitle=\bfseries, title=#1,
  sharp corners, boxrule=0.6pt, breakable
}
\newtcolorbox{noveltybox}[1][]{
  colback=blue!4, colframe=blue!50,
  fonttitle=\bfseries, title=#1,
  sharp corners, boxrule=0.8pt, breakable
}
\newtcolorbox{resourcebox}[1][]{
  colback=green!4, colframe=green!50,
  fonttitle=\bfseries, title=#1,
  sharp corners, boxrule=0.6pt, breakable
}
\newtcolorbox{resultbox}[1][]{
  colback=orange!6, colframe=orange!60!black,
  fonttitle=\bfseries, title=#1,
  sharp corners, boxrule=0.8pt, breakable
}
\newtcolorbox{diagnosticbox}[1][]{
  colback=red!4, colframe=red!50,
  fonttitle=\bfseries, title=#1,
  sharp corners, boxrule=0.8pt, breakable
}

% --- Configuration des Titres ---
\titleformat{\section}
  {\Large\bfseries\color{black}}{\thesection}{1em}{}[\titlerule]
\titleformat{\subsection}
  {\large\bfseries\color{black}}{\thesubsection}{1em}{}
\titleformat{\subsubsection}
  {\bfseries\color{black}}{\thesubsubsection}{1em}{}

% --- Paramètres de paragraphe ---
\setlength{\parskip}{0.5em}
\setlength{\parindent}{0pt}

\title{Interprétabilité automatique de mails clients par Sparse Autoencoders\\
\large Application à la correspondance client d'EDF}
\author{Grégoire Pelletier}
\date{}

% =============================================================================
\begin{document}

% ─────────────────────────────────────────────────────────────────────────────
% PAGE DE GARDE
% ─────────────────────────────────────────────────────────────────────────────
\begin{titlepage}
\centering
\vspace*{1cm}
{\large Université Paris-Saclay --- Faculté des Sciences d'Orsay\par}
{\large Master 2 Mathématiques et Intelligence Artificielle\par}
\vspace{2cm}
{\Huge\bfseries Rapport de stage\par}
\vspace{0.5cm}
{\LARGE Interprétabilité automatique de mails clients par\\ Sparse Autoencoders\par}
\vspace{0.3cm}
{\Large\itshape Application à l'analyse interprétable de la\\ correspondance client d'EDF\par}
\vspace{0.6cm}
{\normalsize\itshape Version soumise au jury --- Master 2 Mathématiques et IA, Université Paris-Saclay\par}
\vspace{1.9cm}

\begin{tabular}{ll}
\textbf{Auteur} & Grégoire Pelletier \\[0.3em]
\textbf{Entreprise d'accueil} & EDF R\&D --- Département SEQUOIA, groupe SOAD \\[0.3em]
\textbf{Lieu} & EDF Lab Paris-Saclay, 7 boulevard Gaspard Monge, 91120 Palaiseau \\[0.3em]
\textbf{Maître de stage} & Philippe Suignard (département SEQUOIA) \\[0.3em]
\textbf{Durée du stage} & 6 mois --- mars 2026 à septembre 2026 \\[0.3em]
\end{tabular}

\vfill
{\large \today\par}
\end{titlepage}

\newpage
\section*{Remerciements}
\addcontentsline{toc}{section}{Remerciements}

Je tiens à remercier Philippe Suignard, Claire Del Gatto et l'ensemble de l'équipe du groupe
Statistiques et Outils d'Aide à la Décision (SOAD) du département SEQUOIA
d'EDF R\&D pour leur accueil, leur disponibilité et les échanges réguliers qui
ont accompagné ce stage.

\newpage
\section*{Résumé}
\addcontentsline{toc}{section}{Résumé}

Ce stage porte sur l'explicabilité automatique des représentations denses de
texte --- en particulier les embeddings, aujourd'hui largement utilisés pour
la classification et la recherche de documents chez EDF, mais dont la
production reste, par construction, une boîte noire. Le cas d'usage retenu
pour ancrer cette investigation est l'analyse de la correspondance client
d'EDF : un volume important de mails dont le traitement automatisé
(priorisation, orientation vers le bon service, détection des réclamations et
des urgences) doit, dans un contexte réglementé, rester justifiable a
posteriori. Les Sparse Autoencoders (SAE) projettent l'espace latent d'un modèle dense et "polysémentique" dans un espace de plus grande dimension. L'objectif est d'obtenir une représentation parcimonieuse dans laquelle chaque direction correspond dans
l'idéal à un concept isolé et nommable en langage naturel. Ils sont étudiés ici
comme réponse au besoin d'explicabilité a posteriori.

Deux pipelines indépendants sont développés. Le premier étend un SAE
préentraîné à grande échelle (GemmaScope, sur les activations d'un modèle de
la famille Gemma-3 --- dont la taille est elle-même un des paramètres
étudiés) par un second SAE entraîné spécifiquement sur le domaine cible,
selon une architecture à \emph{cœur gelé} : le dictionnaire préentraîné reste
figé et seule l'extension s'entraîne, sur le résidu de reconstruction non
expliqué par le cœur, ce qui préserve les capacités générales du SAE
d'origine tout en évitant que les concepts génériques et les concepts du
domaine ne se disputent le même budget de parcimonie. Le second pipeline,
indépendant du premier, ne s'appuie pas sur les activations internes d'un
LLM mais sur des embeddings de phrase produits par des modèles multilingues
dédiés (F2LLM-v2, bge-m3) ; un SAE y est entraîné \emph{from-scratch},
directement sur ces embeddings, sans cœur préentraîné à geler puisqu'aucun
dictionnaire de référence n'existe pour cet espace.

Le corpus combine des mails clients et des variantes augmentées couvrant
plusieurs axes de perturbation (émotion, registre, urgence...), construites
pour représenter la diversité réelle de la correspondance traitée.
L'interprétabilité de chaque feature du dictionnaire est mesurée par un
protocole d'auto-interprétation (\emph{odd-one-out}) : un juge LLM doit
isoler, parmi des documents qui activent fortement une même feature, un
document-contrôle n'ayant aucune raison de partager le concept recherché ; il
n'y parvient que si les autres documents partagent effectivement un concept
identifiable. Les features soumises au juge sont échantillonnées de façon
\textbf{stratifiée par fréquence d'activation} pour
représenter l'ensemble du dictionnaire ; le juge appartient par ailleurs à une \textbf{famille de modèle
différente} de celle du modèle dont les features sont extraites, pour écarter
un biais d'auto-préférence.

Une campagne d'ablations portant sur les hyperparamètres du SAE (largeur,
capacité et parcimonie de l'extension, volume de tokens, graine
d'entraînement, dimension d'embedding, backbone du second pipeline) ne révèle,
sous ce protocole, \textbf{aucun effet mesurable} de ces leviers sur le taux
d'interprétabilité. Sous la configuration retenue comme référence pour ce
rapport, le taux d'interprétabilité atteint \textbf{84,7\%}. Des tests complémentaires --- fidélité et plausibilité
de l'explication au niveau document, robustesse du protocole de jugement,
biais multilingue, fidélité du steering, évaluation quantitative du retrieval
--- achèvent la validation du système.

\textbf{Mots-clés} : Sparse Autoencoders, interprétabilité mécaniste,
GemmaScope, grands modèles de langage, explicabilité, traitement automatique
des mails clients, auto-interprétation par juge LLM.

\section*{Abstract}
\addcontentsline{toc}{section}{Abstract}

This internship addresses the automatic explainability of dense text
representations --- embeddings in particular, now widely used at EDF for
classification and document retrieval, yet a black box by construction. The
use case chosen to ground this investigation is the analysis of EDF's
customer correspondence: a large volume of emails whose automated processing
(prioritization, routing, complaint and urgency detection) must, in a
regulated setting, remain justifiable after the fact. Sparse Autoencoders
(SAE), which relearn a sparse representation in which each direction ideally
corresponds to a single, nameable concept, are studied here as an answer to
that need.

Two independent pipelines are developed. The first extends a large pretrained
SAE (GemmaScope, on the activations of a model from the Gemma-3 family ---
whose size is itself one of the parameters studied) with a second SAE trained
specifically on the target domain, using a \emph{frozen-core} architecture:
the pretrained dictionary stays fixed and only the extension is trained, on
the reconstruction residual left unexplained by the core, which preserves the
original SAE's general capabilities while keeping generic and domain concepts
from competing for the same sparsity budget. The second pipeline, independent
of the first, no longer relies on a LLM's internal activations but on
sentence embeddings produced by dedicated multilingual models (F2LLM-v2,
bge-m3); a SAE is trained \emph{from scratch} directly on these embeddings,
with no pretrained core to freeze since no reference dictionary exists for
that space.

The corpus combines customer emails with augmented variants spanning several
perturbation axes (emotion, register, urgency...), built to represent the
real diversity of the correspondence being processed. Each dictionary
feature's interpretability is measured through an auto-interpretation
protocol (\emph{odd-one-out}): an LLM judge must single out, among documents
that strongly activate the same feature, one control document with no reason
to share the target concept; it can only succeed if the other documents do
share an identifiable concept. Features submitted to the judge are sampled
\textbf{stratified by activation frequency} rather than by magnitude, so as to
represent the whole dictionary rather than only its densest directions; the
judge also belongs to a \textbf{different model family} than the model whose
features are being extracted, to rule out a self-preference bias.

An ablation campaign over the SAE's hyperparameters (width, extension
capacity and sparsity, token volume, training seed, embedding dimension,
Pipeline-2 backbone) reveals, under this protocol, \textbf{no measurable
effect} of any of these levers on the interpretability rate. Under the
configuration adopted as this report's reference, the interpretability rate
reaches \textbf{84.7\%}. The \textbf{extractor/judge model's scale}, identified
in earlier campaigns as the only lever producing a clean, significant effect,
does not replicate at any pair (1B/4B/12B, $p\geq0.29$) once the three scales
are measured under a fully homogeneous protocol (stratified selection,
Qwen3.8-27B judge, deduplication) --- one measurement point (27B) remains in
progress at the time of writing to determine whether this effect, so far the
study's most robust, survives this control. Complementary tests --- document-level explanation
fidelity and plausibility, judge-protocol robustness, multilingual bias,
steering fidelity, quantitative retrieval evaluation --- complete the
system's validation.

\textbf{Keywords}: Sparse Autoencoders, mechanistic interpretability,
GemmaScope, large language models, explainability, customer email analysis,
LLM auto-interpretation.

\newpage
\tableofcontents
\newpage
\listoffigures
\listoftables
\newpage

\section*{Glossaire}
\addcontentsline{toc}{section}{Glossaire}
\begin{description}[leftmargin=3.2cm, style=nextline]
  \item[SAE] Sparse Autoencoder (autoencodeur creux) : réseau qui reconstruit un
    vecteur d'entrée à partir d'une représentation intermédiaire parcimonieuse.
  \item[Feature] Une coordonnée du code latent produit par un SAE. Dans
    l'idéal, chaque feature s'active pour un unique concept sémantique
    identifiable (\emph{ex.} \og{}réclamation liée à la facturation\fg{}), et
    seulement pour lui --- c'est cette correspondance, quand elle est
    vérifiée, qui rend le modèle interprétable.
  \item[LLM] Large Language Model (grand modèle de langage).
  \item[GemmaScope] Famille de SAE préentraînés par Google DeepMind sur les
    activations des modèles Gemma/Gemma-2/Gemma-3 \cite{lieberum2024gemmascope}.
  \item[Neuronpedia] Plateforme externe hébergeant un catalogue de labels
    (générés par LLM, parfois validés humainement) déjà associés à chaque
    feature de nombreux SAE publiés, dont GemmaScope --- utilisée ici comme
    source de labels initiaux, complétée par le juge LLM local de ce stage.
  \item[L0] Pseudo-norme comptant le nombre de coordonnées non nulles d'un
    vecteur --- mesure la parcimonie effective d'un code SAE.
  \item[NMSE] Normalized Mean Squared Error --- erreur de reconstruction
    normalisée par la variance du signal d'origine.
  \item[FVE] Fraction de Variance Expliquée (parfois notée $R^2$).
  \item[MRL] Matryoshka Representation Learning --- apprentissage d'embeddings
    dont les préfixes de dimensions sont eux-mêmes des embeddings valides.
  \item[BM25] Fonction de pondération classique en recherche d'information,
    ici appliquée au vocabulaire latent d'un SAE (Latent Terms).
  \item[odd-one-out] Protocole d'auto-interprétation : un juge LLM doit
    identifier l'intrus parmi des extraits activant fortement une même feature.
  \item[McNemar] Test statistique apparié pour comparer deux proportions
    mesurées sur les mêmes unités d'observation sous deux conditions.
  \item[Cochran-Armitage] Test de tendance pour une variable binaire face à un
    facteur ordonné à plusieurs niveaux.
\end{description}

\newpage
% ─────────────────────────────────────────────────────────────────────────────
\section{Introduction et contexte}
% ─────────────────────────────────────────────────────────────────────────────

\subsection{Contexte industriel}

Ce stage s'est déroulé au sein d'EDF R\&D, dans le département SEQUOIA
(\textbf{S}ervices, \textbf{É}conomie, \textbf{Q}uestions h\textbf{U}maines,
\textbf{O}utils \textbf{I}nnovants et I\textbf{A}), qui éclaire le Groupe EDF
sur les sujets de transitions énergétique, numérique et sociétale du secteur et conçoit
les outils --- data analytics, IA, sciences humaines et sociales --- au
service de la connaissance client et de la relation client. Le stage a été
réalisé au sein du groupe de recherche SOAD (Statistiques et Outils d'Aide à
la Décision), équipe E7S, sur le site d'EDF Lab Paris-Saclay. La mission
initiale, formulée par l'offre de stage
(\emph{cf.}~annexe des sources), porte sur l'explicabilité des
\textbf{embeddings de texte} : des vecteurs denses obtenus par transformation
d'un document textuel, aujourd'hui massivement utilisés (classification,
recherche de documents similaires) mais dont la production reste, par
construction, une boîte noire. Les Sparse Autoencoders (SAE), apparus
récemment en interprétabilité mécaniste, sont identifiés par l'offre de stage
comme une piste prometteuse pour lever ce verrou.

Le cas d'usage retenu pour ancrer cette investigation méthodologique est
l'analyse de mails clients reçus par EDF : un volume important de
correspondance dont le traitement (priorisation, orientation vers le bon
service, détection des réclamations et des urgences) est aujourd'hui coûteux à
automatiser de façon à la fois efficace et interprétable --- un préalable
important dans un contexte réglementé où une décision automatisée doit
pouvoir être justifiée.

\subsection{Fondements théoriques : polysémanticité et superposition}

Les grands modèles de langage offrent une capacité de compréhension fine du
texte, mais leurs représentations internes restent largement opaques : une
même dimension de leur espace d'activation encode typiquement plusieurs
concepts sémantiques non corrélés, un phénomène appelé
\textbf{polysémanticité}. Il découle de l'hypothèse de \textbf{superposition}
\cite{elhage2022superposition} : contraints par une dimensionnalité fixe, les
réseaux de neurones compressent davantage de caractéristiques (\emph{features})
que de dimensions disponibles, en exploitant des directions quasi-orthogonales
de leur espace latent.

Cette question --- dans quelle mesure un LLM « comprend » réellement le texte
qu'il traite plutôt que d'en imiter la surface statistique --- est elle-même
débattue ; Beckmann \& Queloz \cite{beckmannqueloz2026} soutiennent que les avancées récentes de
l'interprétabilité mécanique rendent la position purement sceptique de moins
en moins tenable. Les Sparse Autoencoders proposent précisément de réapprendre
une représentation parcimonieuse et de plus haute dimension dans laquelle
chaque direction --- désormais appelée \textbf{feature} --- correspond, dans
l'idéal, à un concept isolé et nommable en langage naturel
\cite{bricken2023monosemanticity, cunningham2023sparse}. C'est cette
correspondance entre feature et concept que ce rapport cherche à mesurer,
via le taux d'\textbf{interprétabilité} : la proportion de features pour
lesquelles un concept partagé et identifiable est effectivement retrouvé.

\subsection{Objectifs du stage}

L'ambition du stage est de construire une plateforme d'analyse de
mails, les applications sont larges : l'indexation et la recherche par concept, le clustering
interprétable de documents, la détection d'urgence et d'intention, la
comparaison de corpus (\emph{diffing} : repérer automatiquement les features
dont l'activation moyenne diffère significativement entre deux sous-corpus,
\emph{ex.} les mails \og{}urgents\fg{} contre le reste, pour faire émerger ce
qui les distingue), la visualisation des concepts activés, le retrieval par
propriétés, et l'explication des décisions prises en aval de ces
représentations. Il est apparu que la communauté d'interprétabilité mécanistique était assez active, et il était logique de réutiliser au maximum
l'existant (SAELens, GemmaScope, la littérature sur l'interprétation
d'embeddings documentaires) plutôt que de réimplémenter, et de documenter
systématiquement les écarts avec ces références lorsqu'une réimplémentation
partielle s'avère nécessaire.

Une phase préliminaire de recherche bibliographique (état de l'art détaillé au
chapitre~\ref{sec:etat-art}) a permis de dresser un panorama des architectures
SAE et des modèles d'embeddings candidats, et de formuler plusieurs pistes de
contribution méthodologique (architecture hybride à cœur gelé, pipeline
micro/macro, dimensionnement des données). La suite du stage a consisté à
implémenter, tester empiriquement et, le cas échéant, corriger ces
hypothèses de travail sur des données réelles du domaine.

\begin{definitionbox}[Précision importante sur la nature du corpus]
Le corpus d'emails utilisé (\texttt{Mails.tsv}), désigné dans ce rapport par
commodité comme les mails \og{}originaux\fg{}, n'est \textbf{pas} de la
correspondance client authentique : il s'agit d'un jeu de données déjà
synthétique, produit par un travail antérieur du laboratoire EDF R\&D (non
réalisé pendant ce stage). Les variantes \emph{augmentées} sont, elles,
générées pendant ce stage (Gemma-3-12B-it) à partir de ces mails
\og{}originaux\fg{}. Aucune des deux couches du corpus n'est donc constituée
de données réelles au sens strict --- le terme \og{}original\fg{} désigne
uniquement leur statut de donnée d'entrée par rapport aux variantes qui en
dérivent, pas leur authenticité.
\end{definitionbox}

\subsection{Démarche et plan du rapport}

Le travail
s'est structuré autour d'un même fil directeur, répété à chaque étape :
diagnostiquer par ablation contrôlée plutôt que par intuition, puis valider
statistiquement avant de généraliser. Cette démarche a d'abord permis
d'identifier et de corriger un défaut de conception du corpus d'entraînement
de l'extension (\emph{cf.}~\S\ref{subsec:diagnostic-initial}), puis de la
réappliquer systématiquement à une large campagne d'ablations
d'hyperparamètres, jusqu'à la méthodologie statistique elle-même. Ce rapport suit une organisation thématique et non
chronologique : les choix d'architecture et les résultats sont présentés dans
leur forme aboutie, les étapes intermédiaires n'étant mentionnées que
lorsqu'elles éclairent directement une décision (\emph{cf.} notamment
\S\ref{subsec:diagnostic-initial} et \S\ref{subsec:sae-boost-litt}).

Le chapitre~\ref{sec:etat-art} positionne le projet par rapport à l'état de
l'art. Le chapitre~\ref{sec:contribution} décrit l'architecture et la
méthodologie mises en œuvre. Le chapitre~\ref{sec:resultats} présente
l'ensemble des résultats expérimentaux, cœur scientifique du rapport. Le
chapitre~\ref{sec:discussion} discute les limites, les perspectives et conclut.

\newpage
% ─────────────────────────────────────────────────────────────────────────────
\section{Méthodes, algorithmes et état de l'art}
\label{sec:etat-art}
% ─────────────────────────────────────────────────────────────────────────────

\subsection{Fondements théoriques de l'apprentissage de dictionnaires}

L'hypothèse de la représentation linéaire postule que les modèles de langage
représentent les concepts de haut niveau comme des directions linéaires dans
leur espace d'activation \cite{elhage2022superposition}. Soit une distribution
de données $D_x$ sur un espace d'entrée $\mathcal{X}$ (textes), et une
fonction d'activation interne $h : \mathcal{X} \rightarrow \mathbb{R}^{d_\text{model}}$
extraite d'un modèle. L'objectif de l'apprentissage de dictionnaire
parcimonieux est de trouver une matrice de décodage
$W_\text{dec} \in \mathbb{R}^{d_\text{model} \times d_\text{hidden}}$
($d_\text{hidden} \gg d_\text{model}$) et une fonction d'encodage
$f : \mathbb{R}^{d_\text{model}} \rightarrow \mathbb{R}^{d_\text{hidden}}$
produisant un vecteur d'activation latent $z = f(h(x))$ tel que la
reconstruction $\hat{h}(x) = W_\text{dec}\, z + b_\text{dec} \approx h(x)$ soit
fidèle, tandis que $z$ reste très parcimonieux ($\|z\|_0 \ll d_\text{hidden}$) :
chaque coordonnée $z_j$ est l'activation de la $j$-ième \textbf{feature}
(\emph{cf.}~glossaire), et sa nullité pour la quasi-totalité des entrées est
précisément ce qui rend chaque feature active potentiellement spécifique et
interprétable, plutôt que diffuse sur l'ensemble du corpus.

\begin{definitionbox}[Convention $W_\text{dec}$ --- nota bene]
Dans ce rapport, $W_\text{dec} \in \mathbb{R}^{d_\text{model} \times d_\text{hidden}}$
et chaque \textbf{colonne} $W_\text{dec}[:,j]$ est la direction de la $j$-ième
feature dans l'espace d'origine (convention Anthropic/Bricken 2023
\cite{bricken2023monosemanticity}). \textbf{SAELens \cite{bloom2024saelens}} (v6) stocke en revanche
\texttt{W\_dec} de forme \texttt{[d\_hidden, d\_model]} : la reconstruction est
alors $\hat{x} = \texttt{sae.W\_dec.T} \mathbin{@} z + \texttt{sae.b\_dec}$.
\end{definitionbox}

\subsection{Architectures des SAE}

\subsubsection{Standard SAE (ReLU, pénalité $L_1$)}
\[
  z = \text{ReLU}(W_\text{enc}(x - b_\text{pre}) + b_\text{enc}), \qquad
  \mathcal{L}_\text{Standard} = \|x - \hat{x}\|_2^2 + \lambda \|z\|_1
\]
Le gradient $\lambda\,\text{sgn}(z)$ pousse continûment chaque activation
active vers zéro, quelle que soit son amplitude réelle : $z$ sous-estime donc
systématiquement l'intensité du signal qu'il devrait porter, un biais dit de
\textbf{rétrécissement} (\emph{shrinkage}). Ce même gradient, en poussant
aussi les poids d'entrée vers zéro dès qu'une feature devient rarement
active, rend le Standard SAE particulièrement vulnérable aux latents morts
\cite{cunningham2023sparse}.


\subsubsection{Le problème des latents morts}

Un latent $j$ est dit \emph{mort} si son activation reste nulle sur un
horizon de batches consécutifs. Pour un Standard SAE, le gradient de la
pénalité $L_1$ pousse monotoniquement les poids vers zéro dès que la
pré-activation devient négative sur tout le batch, ce qui sature
irréversiblement la ReLU. Le diagnostic empirique se fait via
$\rho_\text{dead} = \frac{1}{d_\text{hidden}} \sum_j \ind[\max_x z_j(x) = 0]$ ;
un $\rho_\text{dead} > 5\%$ signale une compression excessive. Chaque
architecture propose son propre mécanisme anti-mort (séparation porte/magnitude
pour Gated, perte AuxK pour Top-K/BatchTopK, \emph{pre-act loss} et
initialisation de seuil pour JumpReLU) --- \emph{cf.}~état de l'art complet en
annexe des sources bibliographiques du stage.

\subsubsection{Gated SAE}
Sépare la décision d'activation (\emph{gate}) de sa magnitude
\cite{rajamanoharan2024gated} :
\[
  z = \ind(\pi_\text{gate}(x) > 0) \odot \text{ReLU}(\pi_\text{mag}(x)), \qquad
  \mathcal{L}_\text{Gated} = \|x - \hat{x}\|_2^2 + \lambda \|\text{ReLU}(\pi_\text{gate}(x))\|_1
\]
La pénalité $L_1$ ne porte que sur la porte : les magnitudes actives ne
subissent aucun rétrécissement.

\subsubsection{Top-K et BatchTopK}
Le Top-K contraint exactement $k$ latents actifs par exemple. Le
\textbf{BatchTopK} \cite{gao2024topk, bussmann2025batchtopk}, désormais
recommandé par SAELens comme architecture de référence, relaxe cette
contrainte à l'échelle du batch, avec une perte auxiliaire AuxK pour éviter
les latents morts :
\[
  \mathcal{L}_\text{BatchTopK} = \|x - \hat{x}\|_2^2 + \alpha\,\mathcal{L}_\text{AuxK}(e), \qquad
  \mathcal{L}_\text{AuxK}(e) = \bigl\|e - W_\text{dec}[:,\mathcal{S}]\,\text{TopK}(\pi_{\mathcal{D}_\text{dead}}, k_\text{aux})\bigr\|_2^2
\]

\subsubsection{JumpReLU}
Restaure une parcimonie adaptative sans pénalité $L_1$
\cite{rajamanoharan2024jumprelu} :
\[
  z_i = \pi_i(x) \cdot H(\pi_i(x) - \theta_i)
\]
où $\theta$ est un vecteur de seuils appris et $H$ la fonction de Heaviside.
La discontinuité rend le gradient nul presque partout ; les auteurs
utilisent un estimateur passe-droit (STE) avec approximation par noyau (KDE)
pour mettre à jour $\theta$.

\subsubsection{Matryoshka SAE}
Segmente le dictionnaire en sous-groupes emboîtés
$M = \{m_1 < m_2 < \dots < d_\text{hidden}\}$ et somme les erreurs de
reconstruction à chaque palier \cite{bussmann2025matryoshka} :
\[
  \mathcal{L}_\text{Matryoshka} = \sum_{m \in M} \omega_m \bigl\|x - (W_{\text{dec},\,0:m}\,z_{0:m} + b_\text{dec})\bigr\|_2^2
\]
Cette contrainte force les premières dimensions à encoder l'information la
plus saillante, avec une granularité adaptative à la clé (troncature sans
réentraînement).

\begin{table}[h!]
\centering
\caption{Comparaison analytique des architectures SAE}
\label{tab:sae-archs}
\vspace{0.5em}\small
\begin{tabular}{@{}lllll@{}}
\toprule
\textbf{Architecture} & \textbf{Mécanisme} & \textbf{Pénalité}
  & \textbf{Shrinkage} & \textbf{Statut} \\ \midrule
Standard  & $\text{ReLU}(\pi)$                          & $\lambda\|z\|_1$          & Non  & Baseline \\
Gated     & $\ind(\pi_g>0) \odot \text{ReLU}(\pi_m)$   & $\lambda\|\text{ReLU}(\pi_g)\|_1$ & Oui & Déprécié\\
Top-K     & Troncature $k$ max                          & AuxK                      & Oui  & Déprécié \\
BatchTopK & Troncature $n{\times}k$ max (batch)         & AuxK                      & Oui  & \textbf{Utilisé} \\
JumpReLU  & $\pi \cdot H(\pi - \theta)$                 & $\lambda\|z\|_0$ (STE)    & Oui  & Cœur préentraîné \\
Matryoshka & BatchTopK + nesting                        & $\sum_{m\in M}\|\cdot\|^2$ & Varie & Non implémenté \\ \bottomrule
\end{tabular}
\end{table}

\subsection{Choix du backbone d'embedding du second pipeline}

Le second pipeline (embeddings de phrase, indépendant de GemmaScope) s'appuie
sur \textbf{bge-m3} et la famille \textbf{F2LLM-v2} (80--330M paramètres),
retenus pour leur légèreté et, pour F2LLM-v2, leur support natif du
\textbf{Matryoshka Representation Learning} (MRL) \cite{kusupati2022mrl} : un
entraînement contrastif multi-résolution
($\mathcal{L}_{MRL} = \sum_{d \in \mathcal{D}} \omega_d \cdot
\mathcal{L}_\text{contrastive}(z_{1:d})$) qui rend chaque préfixe de
dimensions de l'embedding directement exploitable, sans réentraînement --- une
propriété testée empiriquement au chapitre~\ref{sec:resultats} et qui a
inspiré, dans la littérature SAE, l'architecture Matryoshka SAE (\emph{cf.}
Table~\ref{tab:sae-archs}). Le pipeline principal s'appuie quant à lui sur
GemmaScope, disponible \emph{prêt à l'emploi} avec un catalogue de labels
Neuronpedia déjà constitué --- un choix motivé par la disponibilité de cet
écosystème plutôt que par un score MTEB, discuté au
chapitre~\ref{sec:discussion}.

\subsection{Le SAE résiduel à cœur gelé : deux généalogies concurrentes}
\label{subsec:sae-boost-litt}

La piste architecturale retenue pour ce stage --- geler un SAE préentraîné et
entraîner une extension sur son résidu de reconstruction --- a été développée
en parallèle par ce projet et, dans la littérature publiée, par
\textbf{SAE Boost} \cite{koriagin2025saeboost} (Koriagin et al., juillet 2025) :

\begin{align*}
  e &= x - \hat{x}_\text{core} & \text{(résidu du SAE gelé)} \\
  z_\text{res} &= \text{BatchTopK}_{k_\text{res}}\bigl(\text{ReLU}(W_\text{enc}^{(\text{res})} x + b_\text{enc}^{(\text{res})})\bigr) \\
  \mathcal{L}_\text{Boost} &= \|e - \hat{e}\|_2^2 / \Var(e)
\end{align*}

Koriagin et al. rapportent des gains de +25\% à +59\% en variance expliquée
sur des domaines cibles (chimie, russe, débats ONU) pour moins de 1\% de
dégradation sur le domaine général.

\begin{definitionbox}[Pourquoi geler le cœur plutôt que réentraîner un SAE unique ?]
Si l'on réentraînait ou affinait un unique SAE partagé sur un mélange de texte
générique et de texte du domaine cible, les concepts des deux origines se
disputeraient le même budget fixe de features actives ($k$ dans BatchTopK).
Le SAE, pour minimiser son erreur de reconstruction, tend alors à
\textbf{réutiliser une même feature pour plusieurs concepts non liés} plutôt
que d'en dédier une par concept --- c'est cette \textbf{interférence de
routage} qui dégrade la monosémanticité, pour le domaine cible comme pour le
domaine général. Geler le cœur et n'entraîner qu'une extension sur le résidu
non reconstruit élimine ce partage \emph{par construction} : aucun budget de
sparsité n'est jamais mis en concurrence entre cœur et extension.
\end{definitionbox}

C'est l'architecture concrètement implémentée et
validée à l'échelle dans ce stage (\emph{cf.}~\S\ref{subsec:frozencore}) ---
et non une architecture Matryoshka pleinement hiérarchique et emboîtée, la
piste initialement envisagée en phase bibliographique préliminaire mais dont
la complexité d'entraînement et le besoin en données n'étaient pas justifiés
par les résultats obtenus (\emph{cf.}~chapitre~\ref{sec:resultats}).

\subsection{Sanity checks : les métriques proxy suffisent-elles ?}

Korznikov et al. \cite{korznikov2025sanitychecks} montrent qu'un SAE dont le
décodeur est figé à une initialisation aléatoire (jamais entraîné) atteint,
sur les métriques proxy usuelles (variance expliquée, AutoInterp, sparse
probing), des performances étonnamment proches d'un SAE pleinement entraîné.
Ce résultat impose un critère de validation supplémentaire : toute
amélioration mesurée doit être évaluée \emph{relativement} à cette baseline
aléatoire gelée, et non uniquement en absolu. Ce protocole a été directement
reproduit dans ce stage (\emph{cf.}~\S\ref{subsec:sanity-check-resultats}).

\subsection{Interprétabilité au niveau documentaire}

L'application des SAE à des embeddings de documents entiers (plutôt qu'à des
activations token-level) constitue une extension récente de la littérature
\cite{jiang2024sae_embeddings, decodedense2025}. Le pipeline documentaire
distingue \textbf{interprétabilité globale} (que représente la feature $j$ sur
l'ensemble du corpus ? --- mesurée par juge LLM, protocole \emph{odd-one-out})
et \textbf{interprétabilité locale} (quelles features sont actives pour un
document donné ? --- mesurée par max-pooling documentaire des activations
token-level). Le protocole \emph{odd-one-out}
\cite{templeton2024scaling} présente au juge 9 extraits activant fortement une
même feature et 1 intrus aléatoire ; le juge doit identifier l'intrus, ce qui
ne peut réussir, par construction, que si les 9 extraits partagent un concept
identifiable.

\newpage
% ─────────────────────────────────────────────────────────────────────────────
\section{Contribution : architecture et méthodologie}
\label{sec:contribution}
% ─────────────────────────────────────────────────────────────────────────────

\subsection{Vue d'ensemble du système}

Le système implémente deux pipelines d'analyse interprétable de mails
clients EDF, partageant la même infrastructure de corpus, de stockage et
d'évaluation.

\begin{figure}[h!]
\centering
\begin{tikzpicture}[
  node distance=0.7cm and 0.8cm,
  box/.style={draw, rounded corners=2pt, minimum width=2.6cm, minimum height=0.8cm,
    align=center, font=\scriptsize, fill=gray!10},
  arrow/.style={-Stealth, thick}
]
  \node[box] (corpus) {Corpus principal\\Mails originaux +\\variantes augmentées};
  \node[box, below left=1.2cm and -0.3cm of corpus] (p1) {Pipeline 1\\Gemma-3-12B-it\\+ GemmaScope};
  \node[box, below right=1.2cm and -0.3cm of corpus] (p2) {Pipeline 2\\F2LLM-v2 / bge-m3\\+ PhraseLevelSAE};
  \node[box, below=0.9cm of p1] (ext) {Extension\\FrozenCoreResidualSAE\\(entraînée sur résidu)};
  \node[box, below=0.9cm of p2] (train) {SAE entraîné\\from-scratch};
  \node[box, below=0.9cm of ext] (pool1) {Max-pool\\documentaire};
  \node[box, below=0.9cm of train] (pool2) {Max-pool\\documentaire};
  \node[box, below=0.9cm of pool1] (lab1) {Labels : Neuronpedia\\+ juge LLM local};
  \node[box, below=0.9cm of pool2] (lab2) {Labels : juge LLM\\local (odd-one-out)};

  \draw[arrow] (corpus) -- (p1);
  \draw[arrow] (corpus) -- (p2);
  \draw[arrow] (p1) -- (ext);
  \draw[arrow] (p2) -- (train);
  \draw[arrow] (ext) -- (pool1);
  \draw[arrow] (train) -- (pool2);
  \draw[arrow] (pool1) -- (lab1);
  \draw[arrow] (pool2) -- (lab2);
\end{tikzpicture}
\caption{Architecture générale des deux pipelines d'analyse}
\label{fig:pipeline}
\end{figure}

\subsection{Architecture \texorpdfstring{$\textit{FrozenCoreResidualSAE}$}{FrozenCoreResidualSAE}}
\label{subsec:frozencore}

Pour aligner les embeddings d'origine (conçus majoritairement sur des corpus
anglophones) avec les spécificités sémantiques et techniques du corpus cible,
sans induire d'oubli catastrophique sur les capacités générales du SAE
préentraîné \cite{lieberum2024gemmascope}, le Pipeline~1 gèle le dictionnaire du cœur GemmaScope
($W_\text{dec}^\text{core}$) et instancie une extension de $D_\text{extra}$
features dédiées, entraînée exclusivement sur le résidu non reconstruit :
\[
  e = x - \hat{x}_\text{core} = x - \Bigl(b_\text{dec}^\text{core} +
  \sum_{j \in \mathcal{A}_\text{core}} z_j^\text{core}\, w_{\text{dec},j}^\text{core}\Bigr)
\]
\[
  z_\text{extra} = \text{BatchTopK}_{K_\text{extra}}\bigl(\text{ReLU}(W_\text{enc}^\text{extra} x + b_\text{enc}^\text{extra})\bigr), \qquad
  \hat{x} = \hat{x}_\text{core} + W_\text{dec}^\text{extra}\, z_\text{extra}
\]
Les poids $W_\text{dec}^\text{extra}$ sont initialisés par PCA sur un
sous-ensemble de calibration en français ; l'absence de biais additionnel
($b_\text{dec}^\text{extra}=0$) garantit que sur un flux hors-domaine
l'extension reste quasi-inactive, préservant la géométrie de base. Cette
architecture correspond exactement à celle décrite pour SAE Boost
(\S\ref{subsec:sae-boost-litt}) : les résultats de ce stage constituent donc,
\emph{de facto}, une réplication indépendante et une extension empirique de
cette architecture sur un nouveau domaine (embeddings emails clients français plutôt qu'activation de LLM) et à une
échelle de modèle non couverte par la publication d'origine
(1B/4B/12B~paramètres, \emph{cf.}~\S\ref{subsec:model-scale}).

\subsection{Métriques d'évaluation}

\subsubsection{Fidélité topologique locale : \texorpdfstring{$\rho_\text{SAE}$}{rho\_SAE}}
Pour vérifier que la projection SAE ne dénature pas la géométrie de l'espace
d'origine, on évalue la corrélation de rang de Spearman entre les matrices de
similarité cosinus calculées avant et après passage par le SAE, sur un
sous-ensemble de $N$ documents échantillonnés :
\[
  \rho_\text{SAE} = 1 - \frac{6\sum_{i=1}^n d_i^2}{n(n^2-1)}, \qquad
  d_i = \text{rang}(M_\text{raw}^{(i)}) - \text{rang}(M_\text{sae}^{(i)}), \qquad
  n = \binom{N}{2}
\]
Un $\rho_\text{SAE} \to 1$ valide empiriquement la conservation de la
topologie de l'espace de base.

\subsubsection{Sonde logistique (probing)}
Pour vérifier que la décomposition parcimonieuse préserve le pouvoir
discriminant, deux régressions logistiques multinomiales sont ajustées en
parallèle : $\text{Acc}_\text{raw}$ sur les embeddings denses d'origine
(borne supérieure oracle) et $\text{Acc}_\text{sae}$ sur les codes SAE. Le
différentiel $\Delta = \text{Acc}_\text{raw} - \text{Acc}_\text{sae}$ mesure
la perte d'information discriminante induite par la parcimonie. Au-delà de la
conservation de performance, cette sonde rend le modèle intrinsèquement
explicable : les coefficients $\beta_j^{(k)}$ les plus élevés isolent
directement les concepts sémantiques qui dictent la décision.

\subsubsection{Auto-interprétation par juge LLM}
Le taux d'interprétabilité rapporté dans tout ce rapport est mesuré par le
protocole \emph{odd-one-out} : pour chaque feature échantillonnée, un juge LLM
reçoit 9 documents activant fortement la feature et 1 document-contrôle
(activation quasi nulle), et doit désigner l'intrus. Un juge qui réussit
démontre l'existence d'un concept partagé identifiable ; l'échec ne prouve pas
nécessairement l'absence de concept (bruit du protocole,
\emph{cf.}~\S\ref{subsec:robustesse}). La campagne d'ablations
d'hyperparamètres du SAE (\S\ref{sec:resultats}) utilise le même modèle comme
extracteur et comme juge, par simplicité ; la valeur de référence retenue
pour ce rapport (84,7\%, \S\ref{subsec:correction-methode}) découple les deux
rôles, un juge de famille différente (Qwen3.8-27B) jugeant les features
extraites par Gemma-3-12B-it, pour écarter tout biais d'auto-préférence.

\subsection{Construction du corpus}

Le corpus principal combine les mails « originaux » (3\,480 documents,
\emph{cf.}~encadré \S\ref{sec:etat-art}) et leurs variantes augmentées
(39\,949 variantes acceptées sur 45\,240 générées, soit un taux de rejet de
qualité de 11,7\% --- \emph{cf.}~\S\ref{subsec:augmentation}), produites selon
13 axes de perturbation contrôlée (émotion, registre, orthographe, urgence).
Le split train/test est effectué au niveau du \textbf{mail d'origine} (et non
de la ligne) pour éviter qu'une variante augmentée d'un mail de test ne fuite
dans le train --- une variante et son mail source partagent un contenu
factuel quasi identique, ce qui biaiserait artificiellement les métriques de
classification en aval si elles se retrouvaient de part et d'autre du split.

\subsubsection{Garde-fou qualité de la génération augmentée}
\label{subsec:augmentation}
Chaque variante générée est soumise à un contrôle automatique (longueur
minimale, ratio de longueur par rapport au mail source dans $[0{,}4\,;\,2{,}5]$,
non-identité stricte, préservation des entités numériques du mail d'origine
sauf pour l'axe orthographe). 11,7\% des générations sont rejetées par ce
garde-fou, sans impact sur les résultats rapportés : les variantes rejetées
ne sont jamais chargées en aval de la génération.

\subsection{Protocole de validation et infrastructure}

L'ensemble des expériences a été exécuté sur un cluster SLURM à plusieurs
partitions GPU (A100/H100), sans accès réseau direct depuis les nœuds de
calcul --- toutes les dépendances (modèles, données) doivent être
prépositionnées sur disque avant soumission. Gemma-3 présente des activations
de norme très élevée (\emph{massive activations}) dans son flux résiduel ; une
précision fp16 (plage dynamique $\sim 65\,504$) provoque un dépassement
silencieux vers \texttt{inf}/\texttt{nan} sur ces valeurs, corrompant tout
entraînement en aval sans erreur explicite --- le projet utilise bf16 par
défaut partout, qui partage la plage d'exposant de fp32.

Le taux d'interprétabilité étant une proportion mesurée sur un échantillon de
$n=150$ features, chaque comparaison entre deux configurations est testée
statistiquement (test $z$ à deux proportions indépendantes par défaut ; test de
McNemar pour les comparaisons appariées sur les mêmes features ;
Cochran-Armitage pour les plans à niveaux ordonnés --- \emph{cf.}~audit
méthodologique \S\ref{subsec:audit-stats}), avec un seuil conventionnel de
signification $|z| > 1{,}96$.

\newpage
% ─────────────────────────────────────────────────────────────────────────────
\section{Résultats expérimentaux}
\label{sec:resultats}
% ─────────────────────────────────────────────────────────────────────────────

\subsection{Diagnostic et correction du problème initial}
\label{subsec:diagnostic-initial}

Le pipeline complet fonctionnait de bout en bout, mais le run disponible en
début de cette phase, jugé sur un échantillon initial d'une dizaine de
features seulement (corpus generic énergie/sport/support) --- bien trop
restreint pour être statistiquement exploitable, mais suffisant pour alerter
--- échouait le test \emph{odd-one-out} dans la grande majorité des cas.
L'inspection qualitative des exemples présentés au juge a révélé l'absence de
concept commun identifiable entre les 9 extraits
positifs d'une même feature (\emph{ex.} : rappel produit iPad, système
carcéral norvégien, recette de cuisine et agriculture canadienne présentés
ensemble). La lecture du code d'assemblage du corpus a montré la cause
racine : le réservoir de résidus servant à entraîner l'extension était
constitué \textbf{exclusivement} de texte générique (Wikipédia filtré par
mots-clés) --- les mails du domaine cible n'entraient jamais dans les données
d'entraînement de l'extension, seulement dans une visualisation post-hoc.

\begin{resultbox}[Correction et validation]
Une fois le corpus principal reconstruit (mails + variantes augmentées,
\S\ref{sec:contribution}) et le nombre de features jugées porté à 150 pour une
puissance statistique correcte, le taux d'interprétabilité passe à
\textbf{45,3\% (68/150)}, à volume de
tokens comparable. Faire varier ce volume d'un facteur 20 (100k à 2M tokens),
à domaine fixé, ne produit \textbf{aucun effet mesurable} (40,7\% / 45,3\% /
44,7\%, écarts non significatifs). \textbf{Le problème initial était donc un
défaut de domaine du corpus, pas un déficit de volume d'entraînement.}
\end{resultbox}

\subsection{Synthèse des ablations d'hyperparamètres du SAE}
\label{subsec:synthese-ablations}

Une fois le domaine du corpus corrigé, une campagne d'ablations systématique
a testé, pour chaque configuration, une modification \textbf{isolée} par
rapport au run principal (référence : 45,3\%, 68/150), toutes comparées par
test $z$ à deux proportions. \textbf{Toute cette campagne, y compris ses trois
lignes significatives, est mesurée sous sélection par magnitude et juge
auto-référent (gemma-3-12b-it) --- les deux biais isolés et corrigés en
\S\ref{subsec:correction-methode}} : les valeurs absolues ci-dessous se lisent
comme un plancher relatif entre configurations, pas comme le taux réel du
dictionnaire (84,7\%, valeur de référence de ce rapport). Ceci concerne en
particulier les deux lignes d'échelle du modèle, requalifiées à la baisse
sous protocole corrigé --- \emph{cf.}~\S\ref{subsec:model-scale}.

\begin{table}[h!]
\centering
\caption{Synthèse des 14 ablations d'hyperparamètres --- taux d'interprétabilité (n=150, sélection par magnitude, juge auto-référent --- lire comme plancher relatif, \S\ref{subsec:correction-methode})}
\label{tab:synthese}
\vspace{0.4em}\small
\begin{tabular}{@{}lccc@{}}
\toprule
\textbf{Configuration} & \textbf{Taux} & \textbf{z vs référence} & \textbf{Signif.} \\ \midrule
Volume 100k tokens & 40,7\% & 0,82 & non \\
Volume 2M tokens & 44,7\% & 0,12 & non \\
Volume 25M tokens (+filler FineWeb2-fr) & 54,0\% & --1,50 & non \\
Largeur SAE core 65k & 43,3\% & 0,35 & non \\
Largeur SAE core 262k & 46,7\% & --0,23 & non \\
Époques $\times$4 (40/100) & 41,3\% & 0,70 & non \\
Capacité combinée ($D{=}2048$, $K{=}64$) & 40,0\% & 0,93 & non \\
$D_\text{extra}=2048$ seul ($K$ fixe) & 46,0\% & --0,12 & non \\
$K_\text{extra}=5$ (optimum du papier SAE Boost) & 54,7\% & --1,62 & non \\
Graine d'entraînement (seed 123) & 47,3\% & --0,35 & non \\
Combiné v12 (65k+époques+capacité, tranche 1-150) & 44,0\% & 0,23 & non \\
\textbf{Sanity check : décodeur figé aléatoire} & \textbf{29,3\%} & \textbf{2,86} & \textbf{oui} \\
\textbf{Échelle du modèle --- 4B} & \textbf{28,0\%} & \textbf{3,12} & \textbf{oui} \\
\textbf{Échelle du modèle --- 1B} & \textbf{12,0\%} & \textbf{6,38} & \textbf{oui} \\
\bottomrule
\end{tabular}
\end{table}

La graine d'entraînement, en particulier, était motivée par la littérature sur l'instabilité des features individuelles d'un SAE \cite{unstablefeatures2026, identifiablesae2026}.

\textbf{Constat central} : sur les 14 configurations testées, \textbf{11
n'ont aucun effet significatif} sur le taux d'interprétabilité --- ni la
largeur du SAE core, ni le volume de tokens, ni la capacité ou la parcimonie
de l'extension (isolées ou combinées), ni la graine d'entraînement, ne
déplacent la mesure au-delà du bruit d'échantillonnage attendu à $n=150$. Les
3 configurations significatives se distinguent nettement des autres, tant en
magnitude qu'en nature : un sanity check architectural et un unique
hyperparamètre, l'\textbf{échelle du modèle}, discuté ci-dessous.

\subsection{Sanity check : le SAE bat-il un décodeur figé aléatoire ?}
\label{subsec:sanity-check-resultats}

Reproduisant le protocole de Korznikov et al.~\cite{korznikov2025sanitychecks}
(\S\ref{sec:etat-art}), un décodeur $W_\text{dec}^\text{extra}$ figé à une
initialisation aléatoire (jamais entraîné) atteint tout de même
\textbf{29,3\%} d'interprétabilité (44/150) --- nettement en dessous du SAE
pleinement entraîné (45,3\%, écart significatif, $z=2,86$) mais loin d'être
nul. Sur la métrique de classification (séparabilité des 14 classes d'axes
de perturbation), l'écart se resserre fortement (93,5\% entraîné contre
91,2\% figé) : \textbf{la classification en aval résiste beaucoup mieux à un
décodeur aléatoire que l'interprétation qualitative}, répliquant partiellement
le constat du papier sur des métriques proxy insuffisantes.

Cette mesure date de la sélection par magnitude et du juge auto-référent ---
rejouée sous méthodologie finale contre R0 (\S\ref{subsec:correction-methode}),
sous les deux schémas d'initialisation aléatoire de Korznikov et al. :
\emph{iso} (gaussienne isotrope normalisée) et \emph{cov} (gaussienne de
covariance réelle, leur schéma \emph{principal} --- empiriquement plus
difficile à battre, Fig.~1 et Tables~2-4 de leur article). Sous \emph{iso},
un décodeur figé aléatoire laisse \textbf{98,9\% des 1024 directions de
l'extension mortes} (contre 5,7\% pour R0) --- l'encodeur n'a mécaniquement
que 11 directions vivantes à proposer au juge stratifié, contre 150 attendues
--- et la fraction de variance expliquée gagnée par l'extension ($\Delta$FVE,
indépendante du juge) s'effondre à \textbf{$+0{,}0001$ contre $+0{,}1393$ pour
R0}, un facteur ${\sim}1400$. Sous \emph{cov}, le schéma que Korznikov et al.
utilisent pour leurs résultats publiés, la baisse est moins extrême mais dans
le même sens : 97,0\% de directions mortes (31 vivantes), $\Delta$FVE
$+0{,}0086$ --- un facteur ${\sim}16$ sous R0. \textbf{Même sous le schéma
d'initialisation le plus favorable au décodeur aléatoire, le SAE entraîné
explique ${\sim}16\times$ plus de variance résiduelle supplémentaire.} Sur
les features vivantes, le taux d'interprétabilité reste plus bas que R0 dans
les deux cas (54,5\%, $n=11$, \emph{iso} ; 80,6\%, $n=31$, \emph{cov}) mais à
un $n$ trop faible pour porter la conclusion seul --- c'est le $\Delta$FVE,
mesuré sans le bruit d'échantillonnage du juge, qui constitue le test de
falsification le plus net de ce stage.

\begin{resultbox}[Biais de construction du négatif \emph{odd-one-out} --- identifié en fin de rédaction]
Une inspection manuelle des exemples présentés au juge (protocole
\S\ref{subsec:model-scale}) a révélé que le document \og{}négatif\fg{}
(censé ne pas partager le concept de la feature, cf.
\S\ref{sec:etat-art}) est, sur R0, \textbf{plus court que chacun des 9
exemples positifs dans 147 features sur 150 (98\%)} --- souvent réduit à un
seul mot de salutation ou de ponctuation (\texttt{<<,>>}, seul, compte pour
40\% des 150 négatifs). Cause : le négatif est choisi par magnitude
d'activation minimale sur un pool de documents quasi-inactivants, mais le
contexte affiché reste centré sur l'argmax \emph{de cette feature sur ce
document} --- un signal de bruit pur qui atterrit de façon disproportionnée
près du début du document. \textbf{Un juge peut résoudre la tâche par
simple longueur/pauvreté de contexte, sans lire le concept partagé par les 9
autres exemples}, pour l'écrasante majorité des 150 features de chaque run
de ce stage --- un biais commun à \emph{toutes} les mesures d'interprétabilité
de ce rapport, puisque la construction du négatif n'avait jamais été modifiée
avant sa découverte. Un correctif (contexte gauche minimal exigé, avec repli
loggé si aucun candidat ne qualifie, \texttt{src/sae/judge.py}) a été
appliqué et R0 est en cours de rejugement à l'identique (même SAE entraîné,
seul le jugement change) pour mesurer l'ampleur réelle de cette contamination
sur le chiffre de référence de ce rapport. \textbf{Ce biais ne réfute pas les
résultats en $\Delta$FVE ci-dessus} (métrique continue, indépendante de la
construction du négatif) \textbf{ni la conclusion d'absence d'effet
d'hyperparamètre/d'échelle} (\S\ref{subsec:model-scale},
\S\ref{subsec:ablations-finales} : un biais partagé par toutes les
configurations ne change pas leur comparaison relative) --- mais il remet en
question la valeur absolue de tout taux d'interprétabilité cité dans ce
document, y compris la référence 84,7\%.
\end{resultbox}

\subsection{Protocole de jugement : sélection des features et choix du juge}
\label{subsec:correction-methode}

Deux choix de protocole, tout aussi susceptibles d'influencer le taux mesuré
qu'un hyperparamètre du SAE, sont eux-mêmes traités comme des variables
d'ablation : la \textbf{méthode de sélection} des features soumises au juge
(\textbf{magnitude} d'activation moyenne --- favorise systématiquement les
features les plus denses, proches de directions génériques --- contre
\textbf{stratifiée} par bins de fréquence, échantillonnage équilibré sur des
strates log-espacées) et le \textbf{choix du juge} (gemma-3-12b-it,
auto-référent car même famille que le modèle extracteur, contre Qwen3.8-27B,
famille indépendante). Le tableau ci-dessous croise ces deux variables à la
configuration retenue comme référence pour l'ensemble de la campagne
d'ablations ($K_\text{extra}=5$, $D_\text{extra}=1024$, 25M tokens, layer 31,
gemma-3-12b-it comme extracteur) :

\begin{table}[h!]
\centering
\caption{Sélection $\times$ juge, à la configuration de référence (25M tokens)}
\label{tab:selection-juge}
\vspace{0.4em}\small
\begin{tabular}{@{}lcc@{}}
\toprule
\textbf{Sélection} & \textbf{Juge gemma-3-12b-it}\footnotemark & \textbf{Juge Qwen3.8-27B} \\ \midrule
Stratifiée & 82,0\% (123/150) & \textbf{84,7\% (127/150) --- référence (R0)} \\
Magnitude & 89,3\% (134/150) --- M1, \S\ref{subsec:ablations-finales} & --- \\
\bottomrule
\end{tabular}
\end{table}
\footnotetext{Le run source du 82,0\% (§82 de \texttt{RESULTS\_TESTS.md}, cache
purgé depuis) ne porte pas de confirmation de déduplication des exemples
positifs par mail parent, contrairement à R0 dont le sidecar la confirme
explicitement (\texttt{doc\_groups\_dedup=true}) --- la colonne pourrait donc
confondre choix du juge et introduction de la déduplication, pas isoler le
juge seul.}

À sélection stratifiée fixée (colonne de gauche vs colonne de droite,
réserve de la note ci-dessus), le passage au juge découplé ne produit
\textbf{aucun écart significatif} à cette configuration
(82,0\%~$\to$~84,7\%, $z=-0{,}62$,
$p=0{,}535$) --- une réplication, à un volume de tokens et un $K_\text{extra}$
désormais cohérents avec le reste de la campagne, de l'observation déjà faite
à plus faible volume (\S\ref{subsec:diagnostic-initial}) : sous sélection
stratifiée, le taux d'interprétabilité ne dépend (quasiment) plus du modèle
qui juge. La case magnitude+gemma-3-12b-it (M1) complètera la comparaison
avec l'effet de méthode de sélection à cette même configuration ; en son
absence, aucun chiffre construit à un volume inférieur au plancher de
crédibilité retenu pour ce rapport (25M tokens) n'est cité ici.

\begin{resultbox}[Valeur de référence retenue]
\textbf{84,7\% (127/150), $K_\text{extra}=5$, $D_\text{extra}=1024$, 25M
tokens, layer 31, sélection stratifiée, juge Qwen3.8-27B (R0) est le taux
d'interprétabilité de référence de ce rapport.} Les chiffres en 45,x\%/2x,x\%
cités dans la campagne d'ablations d'hyperparamètres du SAE
(\S\ref{subsec:synthese-ablations}) restent des comparaisons \emph{internes}
valides --- les deux bras de chaque ablation y sont mesurés sous le même
protocole --- mais leurs valeurs absolues se lisent comme un plancher, pas
comme le taux réel du dictionnaire.
\end{resultbox}

\subsection{L'échelle du modèle : le levier déterminant}
\label{subsec:model-scale}

Quatre échelles du modèle extracteur/juge sont testées sous une configuration
strictement commune (sélection stratifiée, juge Qwen3.8-27B, déduplication
par mail parent, $K_\text{extra}=5$, $D_\text{extra}=1024$, 25M tokens,
corpus mixte), chacune avec son propre GemmaScope dédié et son propre
\emph{layer} à \og{}~2/3 de profondeur\fg{} (1B : layer 13 ; 4B : layer 17 ;
12B : layer 31 ; 27B : layer 40) --- remplace une table antérieure qui
mélangeait, sous une même colonne \og{}Taux interp.\fg{}, des points jugés
par Qwen et d'autres par le modèle extracteur lui-même (biais d'auto-préférence
non contrôlé, \S\ref{subsec:correction-methode}) : les quatre points ci-dessous
sont désormais comparables terme à terme, sous le même protocole de
sélection/jugement --- mais seul le \textbf{triplet} modèle/layer/SAE core
varie d'un point à l'autre, chaque layer étant la meilleure estimation
\og{}~2/3 profondeur\fg{} disponible pour son échelle : la taille du modèle
n'est donc pas isolée du layer dans ce sweep.

\begin{table}[h!]
\centering
\caption{Effet de l'échelle du modèle extracteur, méthodologie totalement homogène (stratifié, juge Qwen3.8-27B, $K_\text{extra}=5$, 25M tokens)}
\label{tab:model-scale}
\vspace{0.4em}\small
\begin{tabular}{@{}lccc@{}}
\toprule
\textbf{Modèle} & \textbf{Layer} & \textbf{Juge} & \textbf{Taux interp.} \\ \midrule
gemma-3-1b-it & 13 & Qwen3.8-27B & 80,0\% (120/150) \\
gemma-3-4b-it & 17 & Qwen3.8-27B & 81,3\% (122/150) \\
gemma-3-12b-it (référence du rapport) & 31 & Qwen3.8-27B & 84,7\% (127/150) \\
gemma-3-27b-it & 40 & Qwen3.8-27B & 85,3\% (128/150) \\
\bottomrule
\end{tabular}
\end{table}

\begin{resultbox}[Résultat central du stage]
Sous protocole totalement homogène, la progression 1B$\to$4B$\to$12B$\to$27B
(80,0\%~$\to$~81,3\%~$\to$~84,7\%~$\to$~85,3\%) \textbf{n'atteint la
significativité à aucune paire} (1B vs 27B : $z=-1{,}17$, $p=0{,}24$ ; 12B vs
27B : $z=-0{,}16$, $p=0{,}87$) ni en tendance (Cochran-Armitage sur les
quatre points). C'est un résultat que ce rapport n'anticipait pas : la
version antérieure de cette section, construite en mélangeant des points
jugés par des modèles différents, décrivait une progression nette jusqu'à
12B puis un plateau --- cette lecture ne survit pas au retrait du biais de
juge. À $n=150$ par bras, ce protocole ne détecte qu'un écart d'au moins
9,6 points ($\alpha=0{,}05$, puissance 80\%, calcul par
\texttt{src/analysis/stats.py::minimum\_detectable\_effect}) : \textbf{les
quatre points ci-dessus sont tous mutuellement compatibles avec l'absence
totale d'effet à cette puissance}, pas seulement entre eux non
significatifs. Sous cette méthodologie, \textbf{aucun des leviers testés
dans ce stage} --- ni les hyperparamètres du SAE
(\S\ref{subsec:synthese-ablations}, \S\ref{subsec:ablations-finales}) ni
l'échelle du modèle extracteur/juge --- ne produit d'effet mesurable sur le
taux d'interprétabilité à cette puissance. Ce constat lui-même reste soumis
à la réserve de \S\ref{subsec:sanity-check-resultats} (biais de construction
du négatif \emph{odd-one-out}, en cours de correction) : les quatre points
partagent la même construction de négatif, donc un biais commun ne changerait
pas leur comparaison relative, mais pourrait déplacer les quatre taux
ensemble.
\end{resultbox}

L'inspection qualitative des hypothèses de diffing générées à chaque échelle
(\emph{cf.}~\S\ref{subsec:diffing}, cohérentes à 12B, confuses à 1B ---
inversion logique cause/conséquence observée) reste, elle, inchangée par ce
qui précède : indépendamment de l'issue chiffrée ci-dessus, la
\textbf{capacité de raisonnement du juge/extracteur} sur des tâches
qualitatives plus exigeantes que l'odd-one-out continue de se dégrader
nettement aux petites échelles.

\subsection{Ablations d'hyperparamètres et de graine sous méthodologie finale}
\label{subsec:ablations-finales}

La campagne d'ablations \S\ref{subsec:synthese-ablations} (14 configurations)
est intégralement mesurée sous sélection par magnitude et juge auto-référent
--- ses valeurs absolues se lisent comme un plancher, pas comme le taux réel
(\S\ref{subsec:correction-methode}). Le tableau ci-dessous reprend les
hyperparamètres les plus consultés sous méthodologie finale (stratifié +
Qwen3.8-27B + déduplication), tous comparés à R0 (\S\ref{subsec:correction-methode},
84,7\%, 127/150) :

\begin{table}[h!]
\centering
\caption{Ablations d'hyperparamètres et de graine contre R0, méthodologie finale}
\label{tab:ablations-finales}
\vspace{0.4em}\small
\begin{tabular}{@{}llccc@{}}
\toprule
\textbf{Run} & \textbf{Écart à R0} & \textbf{Taux interp.} & \textbf{$z$ / $p$ vs R0} & \textbf{$\Delta$FVE} \\ \midrule
R0 & référence & 84,7\% (127/150) & --- & $+0{,}1393$ \\
V1 & seed 123 (bruit run-à-run) & 90,7\% (136/150) & $z=-1{,}58$, $p=0{,}114$ & $+0{,}1399$ \\
V2 & seed 7 (bruit run-à-run) & 82,7\% (124/150) & $z=0{,}47$, $p=0{,}639$ & $+0{,}1402$ \\
A1 & $K_\text{extra}=32$ (défaut historique) & 78,0\% (117/150) & $z=1{,}48$, $p=0{,}138$ & $+0{,}1483$ \\
A2 & $D_\text{extra}=2048$ (capacité $\times$2) & 88,0\% (132/150) & $z=-0{,}84$, $p=0{,}401$ & $+0{,}1468$ \\
A3 & core 65k (défaut : 16k) & 83,3\% (125/150) & $z=0{,}31$, $p=0{,}753$ & $+0{,}1436$ \\
A4 & $\text{batch}_\text{extra}=16384$ (défaut : 1024) & 89,3\% (134/150) & $z=-1{,}20$, $p=0{,}230$ & $+0{,}1410$ \\
A6 & $\text{épochs}_\text{extra}=40$ (défaut : 10) & 88,0\% (132/150) & $z=-0{,}84$, $p=0{,}401$ & $+0{,}1398$ \\
\bottomrule
\end{tabular}
\end{table}

\begin{resultbox}[Le plancher de bruit run-à-run absorbe les écarts d'hyperparamètres]
V1 et V2 (deux graines de la configuration de R0, aucun hyperparamètre
changé) donnent \textbf{84,7\% / 90,7\% / 82,7\%} --- un écart de 8 points
entre les deux graines, lui-même statistiquement "significatif" non corrigé
entre V1 et V2 ($z=2{,}04$, $p=0{,}042$) alors qu'aucun hyperparamètre ne
diffère. \textbf{A1 et A2 tombent tous deux à l'intérieur de cette fourchette
V1--V2} (82,7--90,7\%) : leurs écarts non significatifs contre R0 ne sont donc
pas seulement "non prouvés", ils sont plus petits que le bruit de graine
mesuré sur ce même protocole. C'est la réponse directe à l'objection attendue
sur l'absence de correction multi-tests du tableau \S\ref{subsec:synthese-ablations}
(\S\ref{subsec:audit-stats}) : la variabilité de graine, jamais mesurée avant
cette campagne sous le cache d'extraction partagé, est du même ordre de
grandeur que les écarts d'hyperparamètres eux-mêmes.

La colonne $\Delta$FVE (fraction de variance expliquée par le SAE complet,
core+extension, moins celle du core seul --- calculée sur un échantillon de
4096 tokens, indépendante du juge LLM) raconte une histoire radicalement
plus stable : $+0{,}1393$ à $+0{,}1483$ sur les huit runs, un écart relatif
de moins de 7\% là où le taux d'interprétabilité varie de 12,7 points
(78,0\,--\,90,7\%). \textbf{La reconstruction ne bouge quasiment pas d'un
run à l'autre ; c'est le jugement \emph{odd-one-out} qui introduit
l'essentiel du bruit observé sur le taux d'interprétabilité.} Cette
dissociation est cohérente avec le sanity check
\S\ref{subsec:sanity-check-resultats} : le décodeur figé aléatoire y montre
$\Delta$FVE $\approx 0{,}0001$, un facteur ${\sim}1400$ sous R0, une preuve de
falsification bien plus nette et bien moins bruitée que l'écart
d'interprétabilité mesuré sur ce même run (limité à $n=11$ features actives).
\end{resultbox}

\subsection{Robustesse du protocole de jugement}
\label{subsec:robustesse}

Deux tests évaluent la sensibilité du protocole \emph{odd-one-out} à des
perturbations de surface qui ne devraient, en principe, pas affecter le
jugement :

\begin{itemize}
  \item \textbf{Réordonnancement} : en répétant la question 5 fois par
    feature avec un ordre de présentation différent, seules 30,7\% des
    features obtiennent une décision unanime. Le taux agrégé bouge peu
    (single-shot 45,3\% contre vote majoritaire 48,7\%, test de McNemar
    $p=0,560$, non significatif) mais \textbf{31,3\% des features changent
    individuellement de statut} selon l'ordre de présentation.
  \item \textbf{Langue du juge} \cite{resck2025multilingual, multilingualconcepts2025} : traduire les exemples et le prompt en
    anglais ne change pas significativement le taux agrégé (français 46,9\%
    contre anglais 45,5\%, McNemar $p=0,894$), mais \textbf{38,6\% des
    features changent de statut} selon la langue --- un taux de bruit
    supérieur à celui du simple réordonnancement.
\end{itemize}

Ces deux résultats convergent vers la même conclusion : le protocole
\emph{odd-one-out} à décision unique est \textbf{générique\-ment bruyant} face
à toute perturbation de surface (ordre, langue), pas spécifiquement biaisé
dans une direction donnée. Un vote majoritaire sur plusieurs répétitions
devrait être préféré à une décision \emph{greedy} unique en usage de
production.

\subsection{Qualité de l'explication document-level}

Deux protocoles complémentaires valident que les features citées comme
explication d'une décision de classification y jouent un rôle causal réel,
et non de simples labels plausibles sans lien :

\begin{itemize}
  \item \textbf{Fidélité par ablation} : mettre à zéro les $K$ features les
    plus contributives à une décision de classification produit une chute de
    probabilité prédite nettement supérieure à l'ablation de $K$ features
    aléatoires ou des $K$ moins contributives --- confirmant que les features
    citées comme explication pilotent réellement la décision.
  \item \textbf{Plausibilité par choix forcé} : un juge LLM, confronté aux
    vraies features actives d'un document et à un jeu de features-leurres,
    préfère significativement les vraies features, plus souvent que le
    hasard.
\end{itemize}

\subsection{Fidélité du steering (\texttt{steer\_and\_decode})}

Le \textbf{steering} consiste à intervenir directement sur l'activation d'une
feature (l'amplifier ou la supprimer) pour en observer l'effet causal sur le
comportement du modèle, plutôt que de se contenter de la corrélation
observationnelle qu'établit l'auto-interprétation. Deux protocoles de mesure
ont été comparés, qui n'évaluent pas la même chose : l'\textbf{ablation en
espace de code} mesure l'effet de l'intervention directement sur le vecteur
latent modifié, sans jamais repasser par une sortie réelle --- un test
idéalisé, isolé du bruit de génération. Le \textbf{round-trip
decode/ré-encode} pousse la vérification plus loin : le vecteur modifié est
d'abord décodé (reconverti vers l'espace d'origine, voire vers du texte), puis
ce résultat est ré-encodé et son activation mesurée à nouveau --- ce protocole
teste si l'intervention \emph{survit} au passage par une sortie effectivement
utilisable, et non plus seulement dans l'espace latent interne du SAE. Le
ratio round-trip / ablation directe quantifie cette survie : proche de 0,
l'intervention est effacée dès qu'elle repasse par une sortie réelle (steering
inutilisable en pratique pour cette intention, quel que soit son effet
apparent en espace de code) ; proche de 1, elle est intégralement préservée ;
au-delà de 1, elle est amplifiée par le passage en sortie réelle.

Ce protocole de steering n'avait jamais été testé au-delà d'une vérification
géométrique superficielle. Un test dédié, supprimant les 10 features les plus
explicatives d'une intention puis appliquant ce round-trip, révèle un
comportement \textbf{très hétérogène selon l'intention} : le round-trip
neutralise quasiment l'intervention pour 2 intentions sur 4 testées (ratio
0,00--0,02$\times$), la préserve pour une troisième (0,90$\times$), et
l'amplifie pour la dernière (1,74$\times$). Le steering n'est donc \textbf{pas}
un mécanisme d'intervention causale fiable et prévisible en sortie réelle à
partir du seul test d'ablation en espace de code : une intervention qui paraît
efficace dans l'espace latent peut s'avérer, selon l'intention visée, sans
effet une fois le résultat effectivement produit.

\subsection{Retrieval quantitatif (Latent Terms, fusion et reranking)}

\textbf{Latent Terms} \cite{clavie2026latentterms} détourne BM25 --- une
fonction de pertinence classique en recherche d'information, conçue pour
classer des documents à partir de mots-clés --- en remplaçant les mots-clés
par des \textbf{indices de features SAE} : chaque document est représenté par
l'ensemble de ses features actives (son \og{}vocabulaire latent\fg{}), une
requête en langage naturel est elle-même encodée puis réduite à ses features
actives, et BM25 classe les documents par recouvrement de features avec la
requête --- exactement comme il classerait des documents par recouvrement de
mots avec une requête textuelle. Fidèle au protocole du papier, ce SAE est
entraîné par pure reconstruction sur des activations \textbf{token} d'un
corpus générique hors domaine (jamais sur les mails eux-mêmes), puis appliqué
aux mails entiers.

Une évaluation quantitative (Precision@10, 4 intentions en paraphrase, contre
une baseline TF-IDF) ne montre \textbf{pas de victoire uniforme} d'une
méthode : Latent Terms domine largement sur l'intention au taux de base le
plus faible (\emph{remboursement}, 10,3\% : P@10 0,90 contre 0,30 pour
TF-IDF), à égalité sur deux intentions, et dominé par TF-IDF sur l'intention
au taux de base le plus élevé (\emph{réclamation}, 54,8\% : P@10 0,50 contre
1,00 --- TF-IDF profite d'un vocabulaire homogène sur cette intention
précise). Fusionner les deux classements par \textbf{RRF} (\emph{Reciprocal
Rank Fusion}) puis reranker le top-50 avec un juge LLM \textbf{domine
uniformément les 4 intentions} (P@10~$\geq$~0,90 partout, MAP 0,922 contre
0,761 pour RRF seul et 0,678 pour TF-IDF seul) --- premier résultat sans
compromis de cette section. Le \emph{Rank-Biased Overlap} entre les
classements TF-IDF et Latent Terms reste très bas (0,05 en moyenne) même
quand leur précision est proche : les deux méthodes retrouvent des documents
largement différents en tête de liste, ce que le reranking exploite plutôt
que d'arbitrer entre deux versions du même classement.

\subsection{Diffing entre sous-corpus}
\label{subsec:diffing}

Le \emph{diffing} --- repérer automatiquement les features dont l'activation
diffère entre deux sous-corpus, pour faire émerger ce qui les distingue
(\emph{cf.}~\S\ref{sec:etat-art}) --- fait partie des objectifs initiaux du
stage mais n'avait, jusqu'ici, été validé que qualitativement. Le protocole
suit celui de la littérature citée : une hypothèse (un concept, formulée en
langage naturel à partir des labels de features) est générée automatiquement
à partir de l'écart de fréquence d'activation entre les deux groupes, puis
\textbf{vérifiée} sur un corpus de documents tenus à l'écart de sa
génération --- le \texttt{verification\_rate} mesure la proportion
d'hypothèses dont l'écart de fréquence se confirme au-delà d'un seuil, le
\texttt{coverage} la proportion de documents cibles couverts par au moins une
hypothèse validée. Faute d'un corpus emails déjà annoté par sous-catégorie
suffisamment large pour ce protocole, la validation ci-dessous porte sur un
sous-corpus générique à deux domaines nettement séparés (énergie \emph{vs.}
sport, FineWeb2-fr) --- un banc d'essai pour le \textbf{mécanisme} de diffing
lui-même, avant application au diffing emails (\emph{ex.} urgents \emph{vs.}
non urgents) qui a motivé l'objectif initial.

Une première mesure, à partir des 10 hypothèses les mieux classées par un
score de corrélation archivé (sans passer par une génération structurée),
valide 4 hypothèses sur 10 (\texttt{verification\_rate}~=~40,0\%, IC95\%
Wilson [16,8\,;\,68,7]\%) et couvre 92,5\% des documents cibles ---
principalement porté par 2 des 4 hypothèses valides, pas une couverture
homogène. Remplacer cette sélection par une \textbf{génération structurée}
(un juge LLM produit directement les hypothèses en JSON, avec un score de
confiance explicite, à partir des 200 features dont l'écart de fréquence est
le plus marqué) fait passer le taux de validation à \textbf{80,0\% (8/10)} ---
les 2 hypothèses invalidées sont aussi celles à la confiance de génération la
plus basse du lot (0,75 et 0,65) --- au prix d'une couverture plus faible
(55,0\%). Cette comparaison reste descriptive ($n=10$ hypothèses dans les
deux cas, pas de test formalisé faute d'un plan commun), mais confirme que le
mécanisme de bout en bout (génération, vérification) fonctionne sans signe de
dégénérescence, et que la méthode de génération des hypothèses affecte le
taux de validation au moins autant que pour la sélection de features à juger
(\S\ref{subsec:correction-methode}) --- porter cette même chaîne sur un
diffing emails reste la suite naturelle de ce résultat.

\subsection{Clustering interprétable de documents}
\label{subsec:clustering}

Le clustering interprétable de documents, second objectif initial resté sans
validation quantitative jusqu'ici, a été exercé de bout en bout sur le
corpus emails : à partir d'une requête en langage naturel (\emph{ex.}
\og{}type de réclamation client\fg{}), un juge LLM génère des mots-clés,
sélectionne les latents de l'extension dont les labels s'y rapportent (union
des \emph{top-k} par mot-clé, 61 features sur les 68 déjà jugées
interprétables), puis regroupe les documents par similarité de Jaccard sur
leur profil d'activation binaire sur ces latents --- une correction de code
a été nécessaire ici, le clustering ciblé utilisant jusque-là une similarité
cosinus sur ce même vecteur binaire, une métrique différente de celle
spécifiée. Chaque cluster obtenu est ensuite étiqueté par un juge LLM, son
accuracy mesurée par réassignation LLM à l'étiquette la plus proche, et sa
\textbf{compacité géométrique} comparée, en espace d'embedding dense
(bge-m3), à des échantillons aléatoires de même taille (\emph{z-conductance} :
négatif si le cluster est plus dense que le hasard) :

\begin{table}[h!]
\centering
\caption{Clustering ciblé sur la requête \og{}type de réclamation client\fg{} (n=300 documents, 4 clusters)}
\label{tab:clustering}
\vspace{0.4em}\small
\begin{tabular}{@{}clcc@{}}
\toprule
\textbf{Cluster} & \textbf{Description (LLM)} & \textbf{Accuracy} & \textbf{$z$-conductance} \\ \midrule
0 (175 doc.) & Résiliation de contrat et litiges de facturation & 61,7\% & $-3{,}12$ \\
1 (25 doc.) & Activation de service et réclamations mixtes & 40,0\% & $-2{,}91$ \\
2 (46 doc.) & Réclamations urgentes pour coupures de courant & 56,5\% & $-4{,}32$ \\
3 (54 doc.) & Demandes administratives mixtes (activation/résiliation) & 5,6\% & $-10{,}06$ \\
\bottomrule
\end{tabular}
\end{table}

Les quatre $z$-conductances sont négatives : les clusters obtenus par
activation SAE sont géométriquement plus compacts, en espace d'embedding
dense, que des échantillons aléatoires de même taille --- une structure
réelle, pas un artefact du hasard. L'accuracy varie en revanche fortement
d'un cluster à l'autre (5,6\% à 61,7\%), cohérent qualitativement avec la
littérature citée sur des clusters SAE à variance d'accuracy plus élevée que
des baselines denses : le cluster le moins bien classé (3) porte lui-même
l'étiquette la moins homogène (\og{}demandes administratives mixtes\fg{}),
une correspondance qui a du sens --- compacité géométrique et cohérence
sémantique jugée par LLM ne mesurent pas la même chose, et peuvent diverger
sur un cluster par nature hétérogène.

\textbf{Limite connue des deux résultats ci-dessus} : une seule requête de
clustering testée, aucune baseline dense/instruction-tuned pour situer
l'accuracy en absolu, nombre de clusters fixé sans sélection empirique
(silhouette, \emph{gap statistic}), et --- pour le diffing --- un seul
tirage de corpus par condition ; ces deux protocoles seront réexécutés à
mesure que de nouveaux runs deviennent disponibles.

\subsection{Choix du backbone d'embedding (Pipeline 2)}

Quatre backbones ont été comparés pour le Pipeline~2 : F2LLM-v2 (80M, 160M,
330M) et bge-m3. \textbf{bge-m3 domine nettement} sur la plupart des
métriques de reconstruction (NMSE $-18,8\%$ par rapport au meilleur F2LLM,
taux de features mortes 4 à 5 fois plus faible, meilleure séparabilité du
corpus de diffing) --- un candidat sérieux pour une suite du projet, non
adopté comme défaut de production dans ce rapport (F2LLM-v2-80M reste utilisé
partout ailleurs) faute d'une comparaison chiffrée équivalente sur le taux
d'interprétabilité lui-même, pas seulement la reconstruction. Ce résultat a
nécessité un correctif de code : le pooling documentaire était câblé en dur
sur le dernier token (adapté à un backbone décodeur causal comme F2LLM),
incorrect pour un encodeur bidirectionnel comme bge-m3 (pooling \texttt{[CLS]}
attendu).

Sur l'\textbf{interprétabilité} elle-même (pas seulement la reconstruction),
mesurée sous le même protocole que Pipeline~1 (sélection stratifiée, juge
Qwen3.8-27B, \S\ref{subsec:correction-methode}), le taux atteint
\textbf{58,7\% (88/150)} --- sensiblement plus bas que la référence
Pipeline~1 (84,7\%), alors que la reconstruction Pipeline~2 est elle-même
bonne et non dégradée. Pipeline~2 fonctionne de bout en bout et reconstruit
bien, mais son niveau d'interprétabilité individuelle des features reste, à
ce stade, inférieur à celui de Pipeline~1.

\subsection{Balayage de la dimension d'embedding (MRL)}

La dimension d'embedding (\texttt{MATRYOSHKA\_DIM}, F2LLM) n'avait jamais été
variée --- un fait notable découvert en creusant la question : pour
F2LLM-80M, la dimension par défaut égale exactement la dimension complète du
modèle, rendant toute \og{}troncature\fg{} antérieure un artefact sans effet
réel. Un balayage sur F2LLM-160M (dimensions 64, 128, 320, 640) montre une
\textbf{dégradation graduelle, pas abrupte} : 10\% des dimensions (64/640)
conservent déjà 94\% de la performance de classification du plein embedding
--- cohérent avec (sans le démontrer formellement, faute de confirmation d'un
entraînement MRL strict côté F2LLM) un comportement de type Matryoshka.

\subsection{Rigueur statistique des comparaisons}
\label{subsec:audit-stats}

Le choix du test statistique suit systématiquement la structure de la
comparaison plutôt qu'un test par défaut unique : un test à deux proportions
indépendantes pour deux groupes de features distincts (la majorité des
ablations de ce chapitre), le test de \textbf{McNemar} lorsque les deux
conditions portent sur les mêmes 150 features (biais multilingue, robustesse
du juge --- $p=0{,}894$ et $p=0{,}560$ respectivement), et un test de
\textbf{tendance de Cochran-Armitage}, plus puissant qu'une série de
comparaisons par paires, pour l'effet dose-réponse de l'échelle du modèle sur
un plan à niveaux ordonnés. Aucune correction pour comparaisons multiples
n'est en revanche appliquée aux $\sim$15 tests d'ablation d'hyperparamètres
du SAE (\S\ref{subsec:synthese-ablations}) --- sans conséquence pratique ici,
le seul résultat significatif l'étant à $p<10^{-9}$, mais une limite à lever
pour toute extension future de ce protocole, sur le modèle de ce qui est déjà
fait pour le diffing par feature (Benjamini-Hochberg).

\newpage
% ─────────────────────────────────────────────────────────────────────────────
\section{Discussion et conclusion}
\label{sec:discussion}
% ─────────────────────────────────────────────────────────────────────────────

\subsection{Bilan par rapport aux objectifs initiaux}

\begin{table}[h!]
\centering
\caption{Bilan par rapport aux étapes de l'offre de stage}
\vspace{0.4em}\small
\begin{tabular}{@{}p{5.5cm}p{8.5cm}@{}}
\toprule
\textbf{Étape prévue} & \textbf{Réalisation} \\ \midrule
Partie bibliographique sur les SAE & Faite --- élargie à un panorama des
  architectures, des embeddings, des lois d'échelle, et de la littérature de
  sanity-check récente (chapitre~\ref{sec:etat-art}). \\
Choix d'une implémentation existante & GemmaScope + SAELens comme cœur
  préentraîné, complété d'une extension développée pendant le stage. \\
Conversion des données en embeddings & Deux pipelines indépendants
  (Gemma-3/GemmaScope, F2LLM/bge-m3). \\
Entraînement d'un SAE & Fait, avec diagnostic et correction d'un défaut de
  corpus majeur en cours de route. \\
Analyses, interprétation, visualisation & Dashboard interactif, juge LLM
  automatisé, retrieval, diffing, clustering interprétable, sondes de
  classification. \\
Rédaction d'un rapport & Ce document. \\
\bottomrule
\end{tabular}
\end{table}

\subsection{Points résolus}

Le défaut de conception du corpus initial a été diagnostiqué et corrigé, avec
un gain d'interprétabilité mesuré et statistiquement établi. L'architecture à
cœur gelé a été validée à l'échelle, avec une reconnaissance explicite de son
équivalence structurelle à SAE Boost. Une campagne d'ablations
particulièrement large a permis d'établir, avec une rigueur statistique
correcte (McNemar, correction de multiplicité identifiée, test de tendance),
qu'aucun hyperparamètre du SAE ne module significativement le résultat une
fois le domaine corrigé --- et que deux leviers de \emph{protocole de mesure},
plutôt que d'architecture, dominent tous les autres effets mesurés :
l'échelle du modèle extracteur/juge, et la méthode de sélection des features
à juger. Isoler le choix du juge lui-même de celui de l'extracteur a permis de
montrer que l'effet d'échelle massif observé n'est pas un biais
d'auto-préférence généralisé, mais concentré sur les features que la
sélection par magnitude sur-échantillonne (\S\ref{subsec:correction-methode}).
Les protocoles de fidélité, plausibilité et robustesse de l'explication ont
tous été mis en œuvre et donnent des résultats positifs ou, à défaut,
honnêtement nuancés (steering). Le diffing et le clustering interprétable,
deux objectifs initiaux restés sans validation quantitative jusqu'à une
version récente de ce rapport, le sont désormais
(\S\ref{subsec:diffing}--\ref{subsec:clustering}) : le clustering, exercé
directement sur le corpus emails, produit des regroupements géométriquement
réels et d'accuracy sémantique hétérogène mais cohérente ; le diffing,
encore validé sur un corpus générique proxy plutôt que sur les emails, montre
que la méthode de génération des hypothèses conditionne fortement le taux de
validation obtenu.

\subsection{Points non résolus et perspectives}

\begin{itemize}
  \item L'architecture Matryoshka pleinement hiérarchique (dictionnaires
    emboîtés à plusieurs niveaux), envisagée en phase bibliographique,
    n'a pas été implémentée --- l'architecture à cœur gelé effectivement
    testée n'a qu'un seul niveau d'extension. Au vu du résultat central de ce
    stage (l'échelle du modèle domine tout effet architectural du SAE), une
    architecture Matryoshka multi-niveaux ne semble pas prioritaire pour la
    suite.
  \item Le découplage juge/extracteur (\S\ref{subsec:correction-methode}) ne
    fait varier que le juge, à extracteur fixe (Gemma-3-12B-it) : isoler
    pleinement la capacité d'\emph{extraction} de celle du \emph{jugement}
    demanderait un plan croisé complet --- modèle léger comme extracteur et
    comme juge, modèle léger comme extracteur mais juge capable, et
    l'inverse --- non mené ici.
  \item Le rejugement Qwen3.8-27B du point 12B de la famille utilisée pour le
    balayage d'échelle (\S\ref{subsec:model-scale}, layer 31,
    $K_\text{extra}=5$) est en cours : l'ampleur exacte de l'effet d'échelle
    sous protocole totalement homogène (juge Qwen aux deux bornes 1B/12B)
    reste, à ce stade, encadrée mais pas définitivement chiffrée.
  \item Le diffing (\S\ref{subsec:diffing}) n'a pour l'instant été validé que
    sur un sous-corpus générique proxy (énergie \emph{vs.} sport) --- porter
    la même chaîne (génération d'hypothèses structurée, vérification sur
    documents tenus à l'écart) sur un diffing emails (\emph{ex.} urgents
    \emph{vs.} non urgents), l'objectif initial qui a motivé cette
    fonctionnalité, reste à faire.
  \item CamemBERTav2 et Qwen3-Embedding, identifiés en phase bibliographique
    comme candidats sérieux, n'ont pas été intégrés faute d'écosystème SAE
    préentraîné équivalent à GemmaScope --- une comparaison chiffrée directe
    reste à faire si l'équipe souhaite objectiver ce choix par rapport au
    score MTEB plutôt que par la disponibilité d'infrastructure.
  \item Plusieurs pistes envisagées en phase préliminaire restent
    inexplorées, faute de temps : un pipeline micro/macro à deux SAE alignés
    par \emph{AlignedUMAP} (une projection 2D qui contraint les positions à
    rester comparables d'un sous-corpus à l'autre, pour visualiser leur
    évolution) ; une modélisation par graphe biparti (documents et features
    comme deux ensembles de nœuds reliés par les activations, pour en
    exploiter la structure de communautés) ; et un protocole d'annotation
    humaine avec calcul du $\kappa$ de Fleiss (mesure d'accord inter-annotateurs
    corrigée du hasard), pour remplacer le juge LLM par une validation humaine
    d'un sous-échantillon.
  \item Aucune comparaison chiffrée directe n'a été menée contre les
    baselines alternatives de SAE Boost (Extended SAE à initialisation
    aléatoire/most-active, SAE Stitching, fine-tuning complet).
  \item L'ablation de volume, sans effet mesurable jusqu'à 25M tokens
    (Table~\ref{tab:synthese}), est étendue à 200M tokens --- la borne haute
    exacte du seuil de convergence revendiqué par SAE Boost --- dans un run en
    cours au moment de la rédaction de ce rapport.
\end{itemize}

\subsection{Conclusion générale}

Ce stage a permis de construire, diagnostiquer et valider empiriquement, sur
un cas d'usage industriel réel (mails clients EDF), une architecture SAE à
cœur gelé pour l'interprétabilité de grands modèles de langage. Le résultat
le plus significatif --- l'effet dominant de l'échelle du modèle
extracteur/juge sur le taux d'interprétabilité mesuré, loin devant tout
hyperparamètre du SAE lui-même --- reformule la question de recherche
initiale : au-delà du choix de l'architecture SAE, c'est la capacité du
modèle de langage sous-jacent à raisonner sur un concept partagé qui
conditionne la qualité de l'interprétation automatique obtenue. Cette
conclusion, obtenue par une démarche d'ablation rigoureuse et statistiquement
fondée, constitue une contribution méthodologique transférable à tout projet
similaire d'interprétabilité appliquée.

\newpage
% =============================================================================
% BIBLIOGRAPHIE
% =============================================================================
\addcontentsline{toc}{section}{Bibliographie}
\begin{thebibliography}{99}\small

\bibitem{elhage2022superposition}
  Elhage, N. et al.
  \textit{Toy Models of Superposition}.
  Transformer Circuits Thread, Anthropic, 2022.

\bibitem{bricken2023monosemanticity}
  Bricken, T. et al.
  \textit{Towards Monosemanticity: Decomposing Language Models With Dictionary Learning}.
  Transformer Circuits Thread, Anthropic, 2023.

\bibitem{cunningham2023sparse}
  Cunningham, H., Ewart, A., Riggs, L., Huben, R., \& Sharkey, L.
  \textit{Sparse Autoencoders Find Highly Interpretable Features in Language Models}.
  ICLR 2024. arXiv:2309.08600.

\bibitem{rajamanoharan2024gated}
  Rajamanoharan, S. et al.
  \textit{Improving Dictionary Learning with Gated Sparse Autoencoders}.
  arXiv:2404.16014, 2024.

\bibitem{gao2024topk}
  Gao, L. et al.
  \textit{Scaling and Evaluating Sparse Autoencoders}.
  arXiv:2406.04093, OpenAI, 2024.

\bibitem{bussmann2025batchtopk}
  Bussmann, B., Nabeshima, N., Karvonen, A., \& Nanda, N.
  \textit{Learning Multi-Level Features with Matryoshka Sparse Autoencoders}.
  ICML 2025. arXiv:2503.17547.

\bibitem{rajamanoharan2024jumprelu}
  Rajamanoharan, S. et al.
  \textit{Jumping Ahead: Improving Reconstruction Fidelity with JumpReLU Sparse Autoencoders}.
  arXiv:2407.14435, DeepMind, 2024.

\bibitem{bussmann2025matryoshka}
  Bussmann, B., Nabeshima, N., Karvonen, A., \& Nanda, N.
  \textit{Learning Multi-Level Features with Matryoshka Sparse Autoencoders}.
  ICML 2025. arXiv:2503.17547.

\bibitem{templeton2024scaling}
  Templeton, A. et al.
  \textit{Scaling Monosemanticity: Extracting Interpretable Features from Claude 3 Sonnet}.
  Anthropic, 2024.

\bibitem{bloom2024saelens}
  Bloom, J., Tigges, C., Duong, A., \& Chanin, D.
  \textit{SAELens}. \url{https://github.com/decoderesearch/SAELens}, 2024.

\bibitem{jiang2024sae_embeddings}
  (Collectif MATS).
  \textit{Interpretable Embeddings with Sparse Autoencoders: A Data Analysis Toolkit}.
  arXiv:2512.10092, 2024.

\bibitem{decodedense2025}
  (Anonyme).
  \textit{Decoding Dense Embeddings: Sparse Autoencoders for Interpreting and Discretizing Dense Retrieval}.
  arXiv:2506.00041, 2025.

\bibitem{kusupati2022mrl}
  Kusupati, A. et al.
  \textit{Matryoshka Representation Learning}.
  NeurIPS 2022.

\bibitem{koriagin2025saeboost}
  Koriagin, N., Aksenov, Y., Laptev, D., Gerasimov, G., Balagansky, N., \& Gavrilov, D.
  \textit{Teach Old SAEs New Domain Tricks with Boosting}.
  arXiv:2507.12990, 2025.

\bibitem{korznikov2025sanitychecks}
  Korznikov, A., Galichin, A., Dontsov, A., Rogov, O., Oseledets, I., \& Tutubalina, E.
  \textit{Sanity Checks for Sparse Autoencoders: Do SAEs Beat Random Baselines?}
  arXiv:2602.14111, 2025.

\bibitem{lieberum2024gemmascope}
  Lieberum, T. et al.
  \textit{Gemma Scope: Open Sparse Autoencoders Everywhere All At Once on Gemma 2}.
  arXiv:2408.05147, Google DeepMind, 2024.

\bibitem{resck2025multilingual}
  Resck, L., Augenstein, I., \& Korhonen, A.
  \textit{Explainability and Interpretability of Multilingual LLMs: A Survey}.
  EMNLP 2025.

\bibitem{clavie2026latentterms}
  Clavié, B., Lee, S., Shakir, A., Kato, M. P.
  \textit{Latent Terms: Dense Retrievers Contain Trivially Extractable BM25-ready Zipfian Vocabularies}.
  arXiv:2605.29384, 2026.

\bibitem{unstablefeatures2026}
  (Anonyme).
  \textit{Unstable Features, Reproducible Subspaces}.
  arXiv:2606.12138, 2026.

\bibitem{identifiablesae2026}
  (Anonyme).
  \textit{Toward Identifiable Sparse Autoencoders}.
  arXiv:2605.31245, 2026.

\bibitem{multilingualconcepts2025}
  (Anonyme).
  \textit{Sparse Autoencoders Can Capture Language-Specific Concepts Across Diverse Languages}.
  arXiv:2507.11230, 2025.

\bibitem{beckmannqueloz2026}
  Beckmann, P. \& Queloz, M.
  \textit{Mechanistic Indicators of Understanding in Large Language Models}.
  2026.

\end{thebibliography}

\end{document}